\clearpage
\SourceLeaf{95}%

\begin{center}
\large A BEAUTIFUL THEOREM OF PROBABILITY THEORY.
\end{center}

\noindent\hspace*{\parindent}If [we set]
\[
\int e^{itu}\varphi t\,dt = \psi u \mathord{.} \sqrt{2\pi},
\]
\noindent with the integral taken from $t=-\infty$ to $t=+\infty$, then
\[
\int e^{-itu}\psi u\,du = \varphi t \mathord{.} \sqrt{2\pi},
\]
\noindent with the integral taken in the same way.

The justification of this theorem is contained in the formula
\[
\Sigma\,\varphi\!\left(t+\frac{2\pi k}{\omega}\right) = \frac{\omega}{2\pi}\Sigma\, e^{-ikt\omega}\int \varphi\theta\mathord{.} d\theta\mathord{.} e^{ik\theta\omega},
\]
\noindent where the integral sign presupposes integration from $\theta=-\infty$ to $\theta=+\infty$ and the $\Sigma$ sign extends over all positive and negative integer values of $k$. One need only let $\omega$ decrease indefinitely in this formula.

\Needspace{14\baselineskip}
\begin{center}
REMARKS ON THE THEOREM OF PROBABILITY THEORY.
\end{center}

\noindent\hspace*{\parindent}\textsc{Gauss} entered this theorem on the cover of a book entitled »Opuscula mathematica, 1799---1813«. There is no exact record of when the theorem was found.

A simple rearrangement gives the following entirely equivalent form of the theorem.
\[
2\pi\varphi(t) = \int^{+\infty}_{-\infty}\!\left(du\int^{+\infty}_{-\infty}\!\cos u(t-x)\mathord{.}\varphi(x)\,dx\right).
\]
